\section{Objetivo general y objetivos específicos}
\label{sec:objetivos}

\subsection{Objetivo general}

Diseñar y validar mediante simulación estocástica un módulo de software
analítico para la estimación de estados en un gemelo digital de chancado
primario minero, aplicando el algoritmo del Filtro de Kalman sobre datos
ruidosos de tipo SCADA, con el fin de optimizar la observabilidad y reducir el
error cuadrático medio en la reconstrucción del volumen de carga interna dentro
de un entorno de simulación controlado.

\subsection{Objetivos específicos}

\begin{enumerate}[label=\textbf{OE\arabic*.},leftmargin=*,itemsep=0.6em]

\item \textbf{Caracterizar estadísticamente} las perturbaciones y el ruido de la
instrumentación industrial del proceso de chancado primario, mediante el
análisis de datos de literatura técnica indexada, para parametrizar las matrices
de covarianza del estimador.

\item \textbf{Modelar matemáticamente} la dinámica simplificada de flujo de
mineral y la tasa de retención dentro de la tolva de recepción del chancador,
utilizando ecuaciones lineales en espacio de estados en tiempo discreto.

\item \textbf{Desarrollar computacionalmente} el \emph{pipeline} algorítmico
recursivo del Filtro de Kalman Lineal, programando de manera óptima las
funciones de actualización temporal y corrección de mediciones ruidosas.

\item \textbf{Validar el desempeño} del módulo de software mediante simulaciones
estocásticas de Monte Carlo, cuantificando la reducción del error cuadrático
medio frente a las lecturas crudas y degradadas del sistema SCADA simulado.

\end{enumerate}

\subsection{Reflexión metodológica de acotamiento}

Los objetivos específicos delimitan de manera estricta las metas del proyecto
bajo la regla de un único verbo en infinitivo por objetivo. Se separa
nítidamente la etapa conceptual de modelación matemática (OE2) de la de
desarrollo e implementación computacional (OE3), y la caracterización de los
parámetros del entorno (OE1) precede a ambas, puesto que las matrices de
covarianza son insumo del modelo y no un resultado de él.

Este ordenamiento lógico asegura que cada objetivo específico genere un
entregable técnico auditable y verificable —parámetros de ruido documentados,
ecuaciones validadas, scripts del \emph{pipeline} y curvas estadísticas de
error—, descartando explícitamente elementos ajenos a la ingeniería analítica de
datos. La correspondencia entre cada objetivo, sus actividades, sus entregables
y sus indicadores de cumplimiento se detalla en la sección~\ref{sec:plan}, y su
articulación con las decisiones metodológicas adoptadas se analiza en la
sección~\ref{sec:coherencia}.
